\chapter*{Samenvatting}
\addcontentsline{toc}{chapter}{Samenvatting}
Klinische voorspellingsmodellen kunnen evidence-based klinische besluitvorming ondersteunen door het ziekterisico van een patiënt in te schatten. Voor een veilig en effectief gebruik ervan is echter meer nodig dan alleen een goede gemiddelde onderscheidingskracht: de voorspelde waarschijnlijkheden moeten nauwkeurig (d.w.z. gekalibreerd) en klinisch bruikbaar zijn, en voldoende robuust om rekening te houden met verschillen tussen centra, meetprocedures en patiëntenpopulaties.

Dit proefschrift combineert toegepast en methodologisch onderzoek om de besluitvorming voor vrouwen met adnexale massa's te verbeteren, met een focus op het Assessment of Different NEoplasias in the adneXa (\gls{ADNEX})-model voor de diagnose van eierstokkanker.

Ten eerste heb ik het bewijsmateriaal uit externe validatiestudies voor \gls{ADNEX} samengevat. Een systematische review en meta-analyse van 47 studies met in totaal meer dan 17.000 tumoren bevestigde een sterke onderscheidende waarde (AUC 0,93 voor goedaardige versus kwaadaardige gezwellen), maar toonde ook aan dat de kwaliteit van de methodologische rapportage te wensen overlaat: 91\% van de validaties vertoonde een hoog risico op vertekening en de kalibratie werd zelden beoordeeld. Ten tweede presteerde \gls{ADNEX} in een rechtstreekse systematische review en meta-analyse van 11 studies (8271 tumoren) consequent beter dan de Risk of Malignancy Index (\gls{RMI}) wat betreft discriminatievermogen en klinisch nut; bij een drempelwaarde van 10\% voor het risico op maligniteit was de kans dat de test bruikbaar zou zijn in een hypothetisch nieuw centrum 96\% voor \gls{ADNEX} en 15\% voor \gls{RMI}. Ten derde heb ik \gls{ADNEX} geüpdatet naar \emph{ADNEX2} met behulp van prospectieve gegevens uit meerdere centra, waaronder ook conservatief behandelde patiënten, waardoor het model beter kan worden ingezet in een tweestappenstrategie die begint met de goedaardige, eenvoudige descriptoren. 


Het methodologische deel van dit proefschrift behandelt bredere uitdagingen op het gebied van waarschijnlijkheidsschatting en modelevaluatie. Ik heb onderzocht waarom random forest-modellen met uitstekende trainingsprestaties toch onstabiele of misleidende waarschijnlijkheden kunnen opleveren, en heb aangetoond hoe random forests lokale waarschijnlijkheidspieken kunnen leren waarvoor gangbare standaardinstellingen (volledig volgroeide bomen) contraproductief kunnen zijn wanneer het doel een goed gekalibreerde risicoschatting is. Vervolgens heb ik geclusterde kalibratiemethoden ontwikkeld die expliciet rekening houden met een multicenterstructuur, waardoor kalibratie niet alleen op basis van het gemiddelde kan worden samengevat, maar ook met voorspellingsintervallen die de verwachte variabiliteit tussen centra kwantificeren. Ten slotte stelde ik een raamwerk voor onzekerheid in individuele risicoschattingen voor en illustreerde ik empirisch dat, zelfs met grote trainingssteekproeven, de onzekerheid ongerustmakend groot is. Dit ondersteunt het idee dat modellen altijd met de nodige voorzichtigheid moeten worden geïnterpreteerd en dat beslissingen moeten worden genomen door deskundige artsen die de risicoschatting in de juiste context kunnen plaatsen. Tot slot introduceer ik een algoritme voor variabele selectie dat het klinisch nut optimaliseert, waarbij wordt gewaarborgd dat de geselecteerde voorspellende variabelen klinisch relevant zijn, rekening houdend met de kosten van het meten ervan.

Dit proefschrift levert een bijdrage aan het vakgebied van klinische voorspellingsmodellen door nieuw bewijs te leveren over de prestaties van \gls{ADNEX} en door de ontwikkeling van \emph{ADNEX2}, en door nieuwe methoden te ontwikkelen voor het opstellen en valideren van voorspellingsmodellen die in een breed scala aan klinische contexten kunnen worden toegepast. De bevindingen hebben belangrijke implicaties voor de ontwikkeling en validatie van klinische voorspellingsmodellen, zodat deze uitgroeien tot betrouwbare en bruikbare hulpmiddelen voor de klinische besluitvorming.